\section{取り尽くし法}\label{節:取り尽くし法}
先ほどは長方形を並べたが,三角形を使っても面積を求められる.放物線$y=x^2$と直線$y=x$に囲まれた部分を考えよう.\,$0\le x\le1$で,その面積を$Q$とする.この部分に三角形を順に入れていく.

最初の三角形の頂点を$A(0,0),B(1,1),C\left(\frac12,\frac14\right)$とする.$AB$の中点を$D\left(\frac12,\frac12\right)$とすれば,$CD=\frac14$である.三角形$ABC$を$CD$で分けると,二つの三角形$ACD,BCD$はいずれも底辺が$CD$,その底辺への高さが$\frac12$となる.したがって,
\[
 \triangle ABC=2\times\frac12\times\frac14\times\frac12=\frac18
\]
と求められる.

\begin{figure}[htbp]
\centering
\begin{tikzpicture}[x=6cm,y=4.5cm,>=Stealth]
 \fill[softteal](0,0)--(1,1)--(0.5,0.25)--cycle;
 \fill[blue!15](0,0)--(0.5,0.25)--(0.25,0.0625)--cycle;
 \fill[blue!15](0.5,0.25)--(1,1)--(0.75,0.5625)--cycle;
 \draw[->](-0.07,0)--(1.12,0)node[right]{$x$};
 \draw[->](0,-0.07)--(0,1.16)node[above]{$y$};
 \draw[very thick,tealblack,domain=0:1,samples=90]plot(\x,{\x*\x});
 \draw[thick](0,0)--(1,1)--(0.5,0.25)--cycle;
 \draw[blue!65!black](0,0)--(0.25,0.0625)--(0.5,0.25)--(0.75,0.5625)--(1,1);
 \node[below left]at(0,0){$A$};
 \node[above right]at(1,1){$B$};
 \node[below right]at(0.5,0.25){$C$};
\end{tikzpicture}
\caption{最初の三角形の両側に,同じ作り方で三角形を加える.}\label{図:取り尽くし}
\end{figure}

残った二つの部分にも,横の区間の中点にある放物線上の点を取り,三角形を入れる.面積の比を調べるため,一般に$a<b$とし,放物線上の二点$(a,a^2),(b,b^2)$を結ぶ弦を考える.弦の式は
\[
 y=(a+b)x-ab
\]
だから,放物線との縦の差は
\[
 (a+b)x-ab-x^2=(x-a)(b-x)
\]
である.中点$x=\frac{a+b}{2}$での差は$\frac{(b-a)^2}{4}$となる.弦の中点と放物線上の中点を結ぶ縦の線分で三角形を二つに分ければ,先ほどと同じ計算で,両端と中点から作る三角形の面積は
\[
 \frac12(b-a)\frac{(b-a)^2}{4}=\frac{(b-a)^3}{8}
\]
である.横幅が半分になると,三角形の面積は$\frac18$倍となる.二つ合わせても$\frac14$倍である.

よって,各段階で新しく加える面積は
\[
 \frac18,\quad\frac18\cdot\frac14,\quad\frac18\left(\frac14\right)^2,\quad\ldots
\]
となる.第$0$段階から第$N$段階までの和は,等比数列の和の公式から
\[
 P_N=\frac18\cdot\cfrac{1-\left(\cfrac14\right)^{N+1}}{1-\cfrac14}
 =\frac16\left\{1-\left(\frac14\right)^{N+1}\right\}
\]
であり,\,$\frac16$に近づく.

ただし,三角形の和の極限を求めただけでは,それが図形全体の面積であるとはまだ言えない.残りの面積が$0$に近づくことも確かめよう.第$N$段階の後に残る各部分の横幅は$h=2^{-(N+1)}$,個数は$2^{N+1}$である.弦と放物線の差は最大でも$\frac{h^2}{4}$である.各部分は,弦と,それを下へ$\frac{h^2}{4}$だけ平行移動した直線との間の平行四辺形に入る.その面積は横幅と縦の差の積$\frac{h^3}{4}$だから,残りの面積は全部合わせて
\[
 0\le Q-P_N\le2^{N+1}\frac{h^3}{4}
 =\frac1{4\cdot4^{N+1}}
\]
となる.右辺は$N$を大きくすればいくらでも小さくできる.したがって$Q=\frac16$と分かる.

このように,求める図形と面積の分かる図形との差をいくらでも小さくし,候補の値と異なる可能性をなくす考え方を\textbf{取り尽くし法}という.アルキメデスは『放物線の求積』で,放物線と弦に囲まれた部分の面積が,最初に入れた三角形の面積の$\frac43$倍になることを示した.上の計算も
\[
 \frac16=\frac43\cdot\frac18
\]
となっている.ここで示した座標・数列・極限による説明は,その考え方を現代の記号で確かめたものであり,原著の証明をそのまま再現したものではない.

さらに,直線$y=x$の下の三角形の面積は$\frac12$なので,放物線$y=x^2$の下の面積は
\[
 \frac12-Q=\frac12-\frac16=\frac13
\]
となる.長方形を使った計算と一致した.同じ面積でも,下から直接長方形を並べる代わりに,面積の分かる三角形との差として求めることができたのである.

